% SEGMENT OLP-0619-S01
% Part: methods
% Chapter: induction
% Section: relations

\documentclass[../../../include/open-logic-section]{subfiles}

\begin{document}

\olfileid{mth}{ind}{rel}

\olsection{தொடர்புகளும் சார்புகளும்}

இயல் எண்கள் அல்லது ஒழுங்கான சொற்கள் போன்ற
பொருள்களின் கணத்தைத் தொகுத்தறிதல் முறையில்
வரையறுத்திருந்தால், அவற்றின் \emph{மீதான தொடர்புகளையும்}
அதே முறையில் வரையறுக்கலாம். எடுத்துக்காட்டாக,
ஓர் ஒழுங்கான சொல்~$t_1$, மற்றோர் ஒழுங்கான
சொல்~$t_2$ இன் பகுதியாக வந்தால், முதல்
சொல் அதன் துணைச்சொல் என்க. இதற்கு
$t_1 \sqsubseteq t_2$ என்ற குறியைப்
பயன்படுத்துவோம். ஒவ்வோர் ஒழுங்கான சொல்லும்
தனக்கே துணைச்சொல்: $t \sqsubseteq t$.
இந்தத் தொடர்பைத் தொகுத்தறிதல் முறையில்
பின்வருமாறு வரையறுக்கலாம்:

% SEGMENT OLP-0619-S02
\begin{defn}
  ஒழுங்கான சொல்~$t_1$,~$t_2$ இன்
  துணைச்சொல் என்ற தொடர்பான $t_1
  \sqsubseteq t_2$ ஐ,~$t_2$ மீது
  தொகுத்தறிதலால் பின்வருமாறு வரையறுக்கிறோம்:
  \begin{enumerate}
  \item $t_2$ ஓர் எழுத்தெனில்,
    $t_1 \sqsubseteq t_2$ எனில், எனில்
    மட்டுமே, $t_1 = t_2$.
  \item $t_2$ என்பது $[s_1 \circ s_2]$
    எனில், $t_1 \sqsubseteq t_2$ எனில்,
    எனில் மட்டுமே, $t_1 = t_2$ அல்லது
    $t_1 \sqsubseteq s_1$ அல்லது
    $t_1 \sqsubseteq s_2$.
  \end{enumerate}
\end{defn}

% SEGMENT OLP-0619-S03
எடுத்துக்காட்டாக, இந்த வரையறையால்
$\mathrm{a} \sqsubseteq [\mathrm{b} \circ
\mathrm{a}]$ என்று தெரிகிறது. ஏனெனில்
(2) இன்படி $\mathrm{a} \sqsubseteq
[\mathrm{b} \circ \mathrm{a}]$ எனில்,
எனில் மட்டுமே, $\mathrm{a} = [\mathrm{b}
\circ \mathrm{a}]$ அல்லது
$\mathrm{a} \sqsubseteq b$ அல்லது
$\mathrm{a} \sqsubseteq \mathrm{a}$.
முதல் இரண்டும் பொய்: $\mathrm{a}$ என்பது
$[\mathrm{b} \circ \mathrm{a}]$ உடன்
ஒன்றல்ல; (1) இன்படி
$\mathrm{a} \sqsubseteq \mathrm{b}$
எனில், எனில் மட்டுமே,
$\mathrm{a} = \mathrm{b}$; இதுவும் பொய்.
ஆனால் மீண்டும் (1) இன்படி
$\mathrm{a} \sqsubseteq \mathrm{a}$
எனில், எனில் மட்டுமே,
$\mathrm{a} = \mathrm{a}$; இது மெய்.

% SEGMENT OLP-0619-S04
இந்த வரையறை சரியாகச் செயல்படுவதற்கு,
நாம் இன்னும் நிறுவாத ஓர் உண்மை தேவை
என்பதைக் கவனிக்க வேண்டும்: ஒவ்வோர் ஒழுங்கான
சொல்~$t$ உம் தனி எழுத்தாக இருக்கும்;
அல்லது $t = [s_1 \circ s_2]$ ஆகும்
\emph{தனித்துத் தீர்மானிக்கப்படும்} ஒழுங்கான
சொற்கள் $s_1$, $s_2$ இருக்கும்.
``தனித்துத் தீர்மானிக்கப்படும்'' என்றால்,
$t = [s_1 \circ s_2]$ எனில்,
$s_1 \neq r_1$ அல்லது $s_2 \neq r_2$
ஆக இருக்கும்படி அது \emph{மேலும்}
$= [r_1 \circ r_2]$ ஆக இருக்காது.
அப்படி இருந்தால் (2) தன்னுடனேயே
முரண்படக்கூடும்: $t_2$ ஐ $[s_1 \circ s_2]$
என்று வாசிக்கையில் $t_1 \sqsubseteq t_2$
கிடைக்கலாம்; $t_2$ ஐ $[r_1 \circ r_2]$
என்று வாசிக்கையில் $t_1 \sqsubseteq t_2$
பொய் எனக் கிடைக்கலாம். இது நடக்காது
என்று நிறுவுவதற்கு முன், இது \emph{நடக்கக்கூடிய}
ஓர் எடுத்துக்காட்டைப் பார்ப்போம்.

% SEGMENT OLP-0619-S05
\begin{defn}
  \emph{அடைப்பில்லாத சொற்களைத்} தொகுத்தறிதல்
  முறையில் பின்வருமாறு வரையறுக்கவும்:
  \begin{enumerate}
  \item ஒவ்வோர் எழுத்தும் அடைப்பில்லாத சொல்.
  \item $s_1$, $s_2$ அடைப்பில்லாத சொற்களெனில்,
    $s_1 \circ s_2$ உம் அடைப்பில்லாத சொல்.
  \item வேறெதுவும் அடைப்பில்லாத சொல் அல்ல.
  \end{enumerate}
\end{defn}

% SEGMENT OLP-0619-S06
எடுத்துக்காட்டாக $\mathrm{a}$,
$\mathrm{b} \circ \mathrm{d}$,
$\mathrm{b} \circ \mathrm{a} \circ \mathrm{b}$
ஆகியவை அடைப்பில்லாத சொற்கள். அவற்றுக்குத்
``துணைச்சொல்'' என்பதை மேலே போல வரையறுத்தால்,
இரண்டாம் கூறு இவ்வாறு அமையும்:
\begin{center}
  $t_2 = s_1 \circ s_2$ எனில்,
  $t_1 \sqsubseteq t_2$ எனில், எனில்
  மட்டுமே, $t_1 = t_2$ அல்லது
  $t_1 \sqsubseteq s_1$ அல்லது
  $t_1 \sqsubseteq s_2$.
\end{center}

இப்போது $\mathrm{b} \circ \mathrm{a} \circ
\mathrm{b}$ என்பது $s_1 \circ s_2$
என்ற வடிவில் பின்வருமாறு உள்ளது:
\begin{align*}
  s_1 & = \mathrm{b} \text{ மற்றும்} & s_2 & = \mathrm{a} \circ \mathrm{b}.
\intertext{$r_1 \circ r_2$ என்ற வடிவிலும் அது உள்ளது; அப்போது}
  r_1 & = \mathrm{b} \circ \mathrm{a} \text{ மற்றும்} & r_2 &= \mathrm{b}.
\end{align*}
இப்போது $\mathrm{a} \circ \mathrm{b}$
என்பது $\mathrm{b} \circ \mathrm{a} \circ
\mathrm{b}$ இன் துணைச்சொலா? முதல் வாசிப்பின்படி
ஆம்; இரண்டாம் வாசிப்பின்படி இல்லை.

% SEGMENT OLP-0619-S07
ஓர் ஒழுங்கான சொல் மற்ற ஒழுங்கான சொற்களிலிருந்து
உருவாகும் விதம் தனித்தது என்ற பண்புக்கு
\emph{தனித்த வாசிப்புத்தன்மை} என்று பெயர்.
இவ்வாறு தொகுத்தறிதல் முறையில் வரையறுக்கப்பட்ட
பொருள்களுக்கு தொடர்புகளை வரையறுப்பது முக்கியம்
என்பதால், இந்தப் பண்பு அவற்றுக்கு உண்டு
என்று நாம் நிறுவ வேண்டும்.

\begin{prop}
  $t$ ஓர் ஒழுங்கான சொல் என்க. அப்போது
  $t$ தனி எழுத்தாக இருக்கும்; அல்லது
  $t = [s_1 \circ s_2]$ ஆகும்படி
  தனித்துத் தீர்மானிக்கப்படும் ஒழுங்கான சொற்கள்
  $s_1$, $s_2$ இருக்கும்.
\end{prop}

% SEGMENT OLP-0619-S08
\begin{proof}
  $t$ தனி எழுத்தெனில் நிபந்தனை நிறைவேறுகிறது.
  ஆகவே $t$ தனி எழுத்தல்ல என்க. அப்போது
  தொகுத்தறிதல் வரையறையின்படி, சில ஒழுங்கான
  சொற்கள் $s_1$,~$s_2$ க்கு $t$ என்பது
  $[s_1 \circ s_2]$ என்ற வடிவில் இருக்க
  வேண்டும். இவை தனித்துத் தீர்மானிக்கப்படுகின்றன
  என்பதை, அதாவது $t = [r_1 \circ r_2]$
  எனில் $s_1 = r_1$ மற்றும் $s_2 = r_2$
  என்பதை, இன்னும் காட்ட வேண்டும்.

  ஒழுங்கான சொற்கள் $s_1$, $s_2$, $r_1$,
  $r_2$ க்கு $t = [s_1 \circ s_2]$
  என்றும் $t = [r_1 \circ r_2]$ என்றும்
  என்க. $s_1 = r_1$ மற்றும்
  $s_2 = r_2$ என்று காட்ட வேண்டும்.
  முதலில் $s_1$, $r_1$ ஒன்றாகவே இருக்க
  வேண்டும்; இல்லையெனில் ஒன்று மற்றொன்றின்
  முறையான தொடக்கப் பகுதியாக இருக்கும்.
  ஆனால் \olref[sti]{prop:initial} இன்படி,
  $s_1$,~$r_1$ இரண்டும் ஒழுங்கான சொற்களாக
  இருக்கும்போது அது முடியாது. $s_1 = r_1$
  எனில் $s_2 = r_2$ உம் தெளிவாகவே
  பின்தொடர்கிறது.
\end{proof}

% SEGMENT OLP-0619-S09
சார்புகளையும் தொகுத்தறிதல் முறையில்
வரையறுக்கலாம். எடுத்துக்காட்டாக, ஓர் ஒழுங்கான
சொல்லின் உள்ளடங்கிய $[\dots]$ அடைப்புகளின்
அதிகபட்ச ஆழத்திற்கு அதை அனுப்பும் சார்பு~$f$
ஐப் பின்வருமாறு வரையறுக்கலாம்:

\begin{defn}
  \ollabel{defn:depth} ஓர் ஒழுங்கான சொல்லின்
  \emph{ஆழம்} $f(t)$ பின்வருமாறு
  தொகுத்தறிதல் முறையில் வரையறுக்கப்படுகிறது:
  \[
  f(t) = \begin{cases}
    0 & \text{ $t$ ஓர் எழுத்தெனில்}\\
    \max(f(s_1), f(s_2)) + 1 & \text{ $t = [s_1 \circ s_2]$ எனில்.}
  \end{cases}
  \]
\end{defn}

% SEGMENT OLP-0619-S10
எடுத்துக்காட்டாக,
\begin{align*}
  f([\mathrm{a} \circ \mathrm{b}]) & =
    \max(f(\mathrm{a}),f(\mathrm{b})) + 1 = \\
  &= \max(0, 0) + 1 = 1, \text{ மற்றும்}\\
  f([[\mathrm{a} \circ \mathrm{b}] \circ \mathrm{c}]) & =
    \max(f([\mathrm{a} \circ \mathrm{b}]), f(\mathrm{c})) + 1 = \\
  & = \max(1,0) + 1 = 2.
\end{align*}

% SEGMENT OLP-0619-S11
இங்கே $s_1$, $s_2$ ஒழுங்கான சொற்கள்
என்று எடுத்துக்கொள்கிறோம்; மேலும் ஒவ்வோர்
ஒழுங்கான சொல்லும் தனி எழுத்தாகவோ
$[s_1 \circ s_2]$ என்ற வடிவிலோ
இருக்கும் என்ற உண்மையைப் பயன்படுத்துகிறோம்.
இந்த வடிவில் ஒரே வழியில்தான் இருக்க முடியும்
என்பதும் மீண்டும் முக்கியம். ஏனென்று காண,
முன்பு வரையறுத்த அடைப்பில்லாத சொற்களை
மீண்டும் பாருங்கள். அவற்றுக்கான ஒத்த
``வரையறை'' இவ்வாறு இருக்கும்:
\[
  g(t) =
  \begin{cases}
    0 & \text{ $t$ ஓர் எழுத்தெனில்}\\
   \max(g(s_1), g(s_2)) + 1 & \text{ $t = s_1 \circ s_2$ எனில்.}
  \end{cases}
\]
இப்போது அடைப்பில்லாத சொல்
$\mathrm{a} \circ \mathrm{b} \circ
\mathrm{c} \circ \mathrm{d}$ ஐக் கருதுக.
இதை ஒன்றுக்கு மேற்பட்ட வழிகளில் வாசிக்கலாம்.
எடுத்துக்காட்டாக, $s_1 \circ s_2$
என்ற வடிவில்:
\begin{align*}
  s_1 & = \mathrm{a} \text{ மற்றும்} &
  s_2 & = \mathrm{b} \circ \mathrm{c} \circ \mathrm{d},
\intertext{அல்லது $r_1 \circ r_2$ என்ற வடிவில்}
  r_1 & = \mathrm{a} \circ b \text{ மற்றும்} &
  r_2 &= \mathrm{c} \circ \mathrm{d}.
\end{align*}
முதல் வாசிப்பின்படி $g$ ஐக் கணக்கிட்டால்:
\begin{align*}
  g(s_1 \circ s_2) & =
  \max(g(\mathrm{a}), g(\mathrm{b} \circ \mathrm{c} \circ \mathrm{d})) + 1 =\\ & =
  \max(0,2) + 1 = 3
  \intertext{மற்ற வாசிப்பின்படி கிடைப்பது}
  g(r_1 \circ r_2) & =
  \max(g(\mathrm{a} \circ \mathrm{b}), g(\mathrm{c} \circ \mathrm{d})) + 1=\\ &
  = \max(1,1) + 1 = 2
\end{align*}
ஆனால் ஒரு சார்பு எப்போதும் தனித்த மதிப்பைத்
தர வேண்டும். ஆகவே $g$ க்கான நம்
``வரையறை'' உண்மையில் ஒரு சார்பையே
வரையறுக்கவில்லை.

% SEGMENT OLP-0619-S12
\begin{prob}
  $l$ என்ற சார்புக்குத் தொகுத்தறிதல்
  வரையறை தருக; இங்கே $l(t)$ என்பது
  ஒழுங்கான சொல்~$t$ இன் குறிகளின்
  எண்ணிக்கை.
\end{prob}

% SEGMENT OLP-0619-S13
\begin{prob}
  ஒழுங்கான சொற்கள் $t$ மீது
  கட்டமைப்புத் தொகுத்தறிதலைப் பயன்படுத்தி
  $f(t) < l(t)$ என்பதை நிறுவுக.
  இங்கே $l(t)$ என்பது~$t$ இன் குறிகளின்
  எண்ணிக்கை; $f(t)$ என்பது
  \olref[mth][ind][rel]{defn:depth} இல்
  வரையறுக்கப்பட்ட $t$ இன் ஆழம்.
\end{prob}

\end{document}
